% ============================================================
%  test-latex.tex — Mini tutorial de LaTeX
%  Los comentarios empiezan con % y no aparecen en el PDF.
%  Compila con:  pdflatex test-latex.tex
% ============================================================

% ---------- PREÁMBULO ----------
% Todo lo que va antes de \begin{document} es el preámbulo:
% declaras el tipo de documento y cargas paquetes.
\documentclass[12pt,a4paper]{article} % tipos: article, report, book, letter, beamer (diapositivas)

\usepackage[utf8]{inputenc}     % codificación (acentos, ñ)
\usepackage[spanish]{babel}    % reglas tipográficas en español
\usepackage{amsmath, amssymb}  % matemáticas: entornos y símbolos extra
\usepackage{xcolor}           % colores de texto
\usepackage{hyperref}         % enlaces clicables y referencias
\usepackage{graphicx}         % insertar imágenes (\includegraphics)
\usepackage{listings}         % bloques de código

\title{Tutorial de \LaTeX{}}
\author{César}
\date{\today} % \today = fecha de hoy; puedes poner texto fijo

% ---------- DOCUMENTO ----------
\begin{document}

\maketitle   % imprime título, autor y fecha

\tableofcontents % índice automático (compila 2 veces para actualizarlo)

\newpage

% ============================================================
\section{Estructura del documento}
% ============================================================
La jerarquía de secciones es:
\begin{verbatim}
\part{}  \chapter{} (solo report/book)  \section{}
\subsection{}  \subsubsection{}  \paragraph{}
\end{verbatim}
Con un asterisco la sección no lleva número ni sale en el índice:
\texttt{\textbackslash section*\{Sin numerar\}}

\subsection{Una subsección}
Los párrafos se separan dejando una línea en blanco.
Los espacios múltiples      se      ignoran, y los saltos de línea
simples también. Para forzar un salto de línea usa \texttt{\textbackslash\textbackslash} \\
y para una página nueva \texttt{\textbackslash newpage}.

% ============================================================
\section{Formato de texto}
% ============================================================
\begin{itemize}
  \item \textbf{Negrita} → \texttt{\textbackslash textbf\{...\}}
  \item \textit{Cursiva} → \texttt{\textbackslash textit\{...\}}
  \item \underline{Subrayado} → \texttt{\textbackslash underline\{...\}}
  \item \texttt{Monoespaciada} → \texttt{\textbackslash texttt\{...\}}
  \item \emph{Énfasis} (cursiva que se anida bien) → \texttt{\textbackslash emph\{...\}}
  \item \textsc{Capitulares} → \texttt{\textbackslash textsc\{...\}}
  \item {\color{red} Texto en color} → \texttt{\{\textbackslash color\{red\} ...\}} (requiere \texttt{xcolor})
\end{itemize}

Tamaños relativos:
{\tiny tiny} {\scriptsize scriptsize} {\footnotesize footnotesize}
{\small small} {\normalsize normalsize} {\large large} {\Large Large}
{\huge huge} {\Huge Huge}

% ============================================================
\section{Listas}
% ============================================================
Lista con viñetas (\texttt{itemize}):
\begin{itemize}
  \item Primer elemento
  \item Segundo elemento
  \begin{itemize}
    \item Se pueden anidar
  \end{itemize}
\end{itemize}

Lista numerada (\texttt{enumerate}):
\begin{enumerate}
  \item Primero
  \item Segundo
  \begin{enumerate}
    \item Sub-punto
  \end{enumerate}
\end{enumerate}

Lista de descripciones (\texttt{description}):
\begin{description}
  \item[LaTeX] Sistema de composición tipográfica.
  \item[pdflatex] Compilador que genera PDF directamente.
\end{description}

% ============================================================
\section{Matemáticas}
% ============================================================
Fórmula \textbf{en línea} con \$...\$: la ecuación $E = mc^2$ queda
dentro del párrafo. También vale \texttt{\textbackslash( ... \textbackslash)}: \(a^2+b^2=c^2\).

Fórmula \textbf{destacada} con \texttt{\textbackslash[ ... \textbackslash]}:
\[
  \int_0^{\infty} e^{-x^2}\,dx = \frac{\sqrt{\pi}}{2}
\]

Con numeración usa el entorno \texttt{equation}:
\begin{equation}
  x = \frac{-b \pm \sqrt{b^2 - 4ac}}{2a}
  \label{eq:cuadratica}
\end{equation}
y la citas así: la ecuación~\eqref{eq:cuadratica} es la fórmula cuadrática
(\texttt{\textbackslash label} + \texttt{\textbackslash eqref} o \texttt{\textbackslash ref}).

\subsection{Piezas habituales}
\begin{itemize}
  \item Fracción: $\frac{a}{b}$ → \texttt{\textbackslash frac\{a\}\{b\}}
  \item Sub/superíndices: $x_1$, $x^2$, $x_1^2$ → \texttt{x\_1}, \texttt{x\^{}2}
  \item Raíz: $\sqrt{x}$, $\sqrt[3]{x}$ → \texttt{\textbackslash sqrt\{x\}}, \texttt{\textbackslash sqrt[3]\{x\}}
  \item Suma y producto: $\sum_{i=1}^{n} i$, $\prod_{k=1}^{n} k$
  \item Integral: $\int_a^b f(x)\,dx$
  \item Límite: $\lim_{x \to 0} \frac{\sin x}{x} = 1$
  \item Letras griegas: $\alpha, \beta, \pi, \Omega$ → \texttt{\textbackslash alpha}, \texttt{\textbackslash Omega}
  \item Relaciones: $\leq, \geq, \neq, \approx, \in, \subset$
  \item Operadores lógicos: $\forall, \exists, \Rightarrow, \Leftrightarrow$
\end{itemize}

Matrices con \texttt{pmatrix} (paréntesis), \texttt{bmatrix} (corchetes):
\[
  A = \begin{pmatrix} a & b \\ c & d \end{pmatrix},
  \quad
  B = \begin{bmatrix} 1 & 0 \\ 0 & 1 \end{bmatrix}
\]

Varias ecuaciones alineadas con \texttt{align*} (\& marca la alineación,
* quita la numeración):
\begin{align*}
  (a+b)^2 &= (a+b)(a+b) \\
          &= a^2 + 2ab + b^2
\end{align*}

% ============================================================
\section{Tablas}
% ============================================================
El entorno \texttt{tabular} define las columnas:
\texttt{l} = izquierda, \texttt{c} = centro, \texttt{r} = derecha,
\texttt{|} = línea vertical, \texttt{\textbackslash hline} = línea horizontal.

\begin{center}
\begin{tabular}{|l|c|r|}
  \hline
  \textbf{Nombre} & \textbf{Tipo} & \textbf{Valor} \\
  \hline
  pi     & constante & 3.1416 \\
  e      & constante & 2.7183 \\
  phi    & constante & 1.6180 \\
  \hline
\end{tabular}
\end{center}

Con \texttt{p\{4cm\}} la columna tiene ancho fijo y ajusta el texto.
El entorno \texttt{table} + \texttt{\textbackslash caption} hace la tabla
``flotante'' y con número (la posiciona \LaTeX{} donde quede mejor).

% ============================================================
\section{Imágenes}
% ============================================================
Necesitas \texttt{graphicx} y el archivo de imagen junto al .tex:
\begin{verbatim}
\begin{figure}[h]        % h = aquí, t = arriba, b = abajo, p = página
  \centering
  \includegraphics[width=0.8\textwidth]{imagen.png}
  \caption{Pie de foto}
  \label{fig:ejemplo}
\end{figure}
% y luego: ver la figura~\ref{fig:ejemplo}
\end{verbatim}

% ============================================================
\section{Enlaces y referencias}
% ============================================================
Con \texttt{hyperref}:
\begin{itemize}
  \item Web: \url{https://www.latex-project.org} → \texttt{\textbackslash url\{...\}}
  \item Con texto: \href{https://www.overleaf.com/learn}{documentación de Overleaf}
        → \texttt{\textbackslash href\{url\}\{texto\}}
  \item Referencia interna: la sección~\ref{sec:codigo} habla de código
        → \texttt{\textbackslash label} + \texttt{\textbackslash ref}
\end{itemize}

% ============================================================
\section{Código fuente}
\label{sec:codigo}
% ============================================================
En línea: \verb|print("hola")| con \texttt{\textbackslash verb|...|}.

Bloque con \texttt{verbatim} (tal cual, sin formato) o \texttt{lstlisting}
(con colores, del paquete \texttt{listings}):
\begin{lstlisting}[language=Python]
def fibonacci(n):
    a, b = 0, 1
    for _ in range(n):
        yield a
        a, b = b, a + b
\end{lstlisting}

% ============================================================
\section{Citas y notas al pie}
% ============================================================
Texto citado largo con \texttt{quote}:
\begin{quote}
  ``LaTeX es el estándar de facto para documentos científicos.''
\end{quote}

Una nota al pie se inserta en medio del texto\footnote{Así se ve una
nota al pie: \texttt{\textbackslash footnote\{...\}}.}.

Para bibliografía real se usa BibTeX/Biber con un archivo \texttt{.bib}
y comandos \texttt{\textbackslash cite\{clave\}}.

% ============================================================
\section{Trucos finales}
% ============================================================
\begin{itemize}
  \item Caracteres especiales, escápalos con \texttt{\textbackslash}:
        \% \$ \& \# \_ \{ \}
  \item Espacio irrompible (no se corta la línea): \texttt{sección\textasciitilde\textbackslash ref\{...\}}
  \item Puntos suspensivos correctos: \ldots{} → \texttt{\textbackslash ldots}
  \item Comillas tipográficas: ``dobles'' y `simples'
  \item Comentar varias líneas rápido: selecciona y usa \texttt{\%} al inicio
        (en la mayoría de editores, Ctrl+/)
  \item Errores de compilación: lee el \texttt{.log}; el error real suele
        estar en la primera línea que menciona.
\end{itemize}

\end{document}
